\chapter*{Abstrak}
\addcontentsline{toc}{chapter}{\MakeUppercase{Abstrak}}

\small Abstrak merupakan ringkasan singkat dari keseluruhan isi tugas akhir yang memuat latar belakang, tujuan, metode penelitian, hasil, dan kesimpulan. Abstrak bertujuan memberikan gambaran umum kepada pembaca mengenai penelitian yang dilakukan sehingga pembaca dapat memahami pokok permasalahan, pendekatan yang digunakan, serta hasil yang diperoleh tanpa harus membaca keseluruhan dokumen.

Abstrak umumnya ditulis dalam satu hingga tiga paragraf dengan panjang yang disesuaikan dengan ketentuan yang berlaku. Penulisan abstrak harus menggunakan bahasa yang ringkas, jelas, dan informatif serta menghindari penggunaan sitasi, tabel, gambar, maupun persamaan.

Secara umum, isi abstrak dapat disusun sebagai berikut:

\begin{itemize}[topsep=0pt,itemsep=0pt,parsep=0pt,leftmargin=*]
\item Paragraf pertama menjelaskan latar belakang permasalahan, tujuan penelitian, serta ruang lingkup atau batasan penelitian.
\item Paragraf kedua menjelaskan metode, pendekatan, algoritma, atau tahapan yang digunakan untuk menyelesaikan permasalahan yang diteliti.
\item Paragraf ketiga menjelaskan hasil utama penelitian, tingkat keberhasilan yang dicapai, serta kesimpulan yang diperoleh. Apabila memungkinkan, hasil penelitian disajikan secara kuantitatif menggunakan parameter atau indikator kinerja yang relevan.
\end{itemize}

Abstrak harus mampu menjawab secara ringkas pertanyaan berikut:

\begin{enumerate}[topsep=0pt,itemsep=0pt,parsep=0pt,leftmargin=*]
\item Apa permasalahan yang diteliti?
\item Mengapa permasalahan tersebut penting untuk diselesaikan?
\item Bagaimana penelitian dilakukan?
\item Apa hasil yang diperoleh?
\item Apa kesimpulan utama dari penelitian?
\end{enumerate}

Setelah abstrak, sertakan beberapa kata kunci (\textit{keywords}) yang mewakili topik utama penelitian untuk memudahkan proses pengindeksan dan pencarian dokumen.
\vspace{1em}

\noindent \textbf{Kata Kunci:} Tuliskan dua sampai enam kata kunci yang berkaitan dengan masalah yang dibahas.
